\subsection{电场做功与磁场偏转的不同作用}
对于电荷量为 $q$、静质量为 $m$ 的粒子，静电场中的能量增量为 $q(\varphi_i-\varphi_f)$。电子从静止经加速电压大小 $U$ 获得 $K=eU$，所以
\[
 \gamma=1+\frac{eU}{mc^2},\quad
 v=c\sqrt{1-\left(1+\frac{eU}{mc^2}\right)^{-2}},\quad
 pc=\sqrt{eU(eU+2mc^2)}.
\]
磁场力与速度垂直，故不改变总能量。若速度垂直于匀强磁场，动量方向以角速度 $v/r$ 转动，$|\dd\vect p/\dd t|=pv/r$；与 $|q|vB$ 比较得
\[
 r=\frac p{|q|B},\qquad \omega=\frac{|q|B}{\gamma m}.
\]
高速时仍可用“动量除以电荷磁场”求半径，但不能再把动量直接写成 $mv$，也不能把回旋频率当作与能量无关的常数。低能极限两式恢复经典结果。
